% 01-dat-van-de.tex ;  Recap BPTT from ch01, state the empirical question
%
% Scope: remind reader of the BPTT recurrence; why the weight reuse matters;
% set up the question: does the RNN actually learn long-range dependencies?

\section{Điều chương 01 không trả lời}
\label{sec:ch02-dat-van-de}

Chương~01 dẫn xuất BPTT đến từng Jacobian. Ở trang cuối của dẫn xuất đó, tôi
viết một câu mà bạn nên đọc lại: ``$\mat{W}_{hh}\transpose$ xuất hiện ở mọi
bước lùi. Nếu bạn tháo cuộn công thức này từ $t = T$ về $t = 1$, bạn sẽ thấy
$\partial L / \partial \vect{h}_1$ chứa $\mat{W}_{hh}\transpose$ nhân với chính
nó $T$ lần.''

Đó là một phát biểu đại số. Nó nói rằng gradient \emph{có thể} suy giảm hoặc
bùng nổ, chứ chưa nói rằng nó \emph{thực sự} suy giảm hay bùng nổ trong một
bài toán cụ thể. Chương này trả lời câu hỏi đó bằng thực nghiệm.

Tôi sẽ huấn luyện một RNN đơn giản; đúng loại RNN mà chương~01 đã dẫn
xuất; trên hai bài toán tổng hợp kinh điển:
\tn{bài toán cộng}{adding problem} và \tn{bài toán sao chép}{copy task}. Cả
hai được thiết kế để bắt buộc mạng phải nhớ thông tin qua nhiều bước thời
gian. Nếu RNN có thể học phụ thuộc dài hạn, nó sẽ giải được cả hai. Nếu
không, ta sẽ thấy loss không giảm và gradient biến mất ở những bước xa đầu
ra; và đó chính là thứ ta sẽ đo.

Tôi chọn hai bài toán này vì chúng cô lập vấn đề. Không có từ vựng, không có
embedding, không có token hóa. Chỉ có những con số và một câu hỏi: mạng có
nhớ được không?

Trước khi chạy thí nghiệm, tôi cần nói rõ một điều. Chương này \textbf{không}
dẫn xuất điều kiện toán học cho vanishing gradient. Đó là việc của chương~3,
dựa trên bài báo của Pascanu và cộng sự~2013. Chương này chỉ \textbf{đo}:
huấn luyện RNN, ghi lại loss, ghi lại norm của gradient tại từng bước thời
gian, và cho bạn thấy những con số đó nói gì.
